\documentclass[11pt]{article}
\usepackage[margin=1in]{geometry}
\usepackage{enumitem}
% =============================================================================
% Preambles/header.tex — shared LaTeX/Beamer preamble for this project
%
% Use from a Beamer lecture by prepending:
%     \input{header}
% and compiling with TEXINPUTS=../Preambles:$TEXINPUTS (the /compile-latex
% skill does this automatically).
%
% PALETTE CONTRACT. Color names defined here MUST match the names in
% Quarto/theme-template.scss so Beamer and Quarto renderings use the same
% palette. The scripts/check-palette-sync.sh script enforces this.
%
% To customize: edit the HEX values below AND the matching $variable in
% Quarto/theme-template.scss, then run scripts/check-palette-sync.sh.
% =============================================================================

% -----------------------------------------------------------------------------
% Engine detection + encoding
%
% Our /compile-latex skill uses XeLaTeX, where inputenc is unnecessary and
% can produce encoding warnings. Only load inputenc under pdfTeX.
% -----------------------------------------------------------------------------
\usepackage{iftex}
\ifPDFTeX
  \usepackage[utf8]{inputenc}
\fi

\usepackage{amsmath, amssymb, amsthm}
\usepackage{booktabs}
\usepackage{graphicx}

% -----------------------------------------------------------------------------
% PALETTE — mirrors Quarto/theme-template.scss variable names.
% Load xcolor and define colors BEFORE anything that references them
% (notably hyperref, hypersetup, and the Beamer color assignments).
% -----------------------------------------------------------------------------
\usepackage{xcolor}

\definecolor{primary-blue}{HTML}{012169}      % headings, accents
\definecolor{primary-gold}{HTML}{B9975B}      % emphasis, borders
\definecolor{highlight-yellow}{HTML}{F2A900}  % markers, alerts
\definecolor{light-bg}{HTML}{E8EDF5}          % title / section backgrounds
\definecolor{jet}{HTML}{1A1A1A}               % body text

% Semantic colors (used in diagrams and annotations)
\definecolor{positive}{HTML}{15803D}          % good / observed / identified
\definecolor{negative}{HTML}{B91C1C}          % bad / problematic / bias
\definecolor{neutral}{HTML}{525252}           % reference / context

% Highlight accents (match .hi-* classes in the SCSS)
\definecolor{hi-slate}{HTML}{314F4F}
\definecolor{hi-green}{HTML}{2E8B57}
\definecolor{hi-red}{HTML}{C41E3A}

% Semantic aliases so skill templates and snippets can write
% \textcolor{accent}{...} without hard-coding "primary-blue".
\colorlet{accent}{primary-blue}
\colorlet{accent2}{primary-gold}

% -----------------------------------------------------------------------------
% Hyperref — loaded AFTER xcolor + palette so link colors resolve.
% -----------------------------------------------------------------------------
\usepackage{hyperref}
\hypersetup{colorlinks=true, linkcolor=primary-blue, urlcolor=primary-blue,
            citecolor=primary-blue, pdfborder={0 0 0}}

% -----------------------------------------------------------------------------
% Course code — set once per deck (after \input{header}, before \begin{document})
% via \coursecode{CS401}. Shown in the footer on every slide so a deck's course
% is identifiable without opening the title page. Empty by default.
% -----------------------------------------------------------------------------
\makeatletter
\newcommand{\coursecode}[1]{\gdef\@coursecodetext{#1}}
\gdef\@coursecodetext{}
\makeatother

% -----------------------------------------------------------------------------
% Beamer theme assignments (only apply under Beamer).
%
% We detect Beamer via \@ifclassloaded{beamer}, which is the documented,
% non-internal way. This must be wrapped in \makeatletter/\makeatother
% because \@ifclassloaded contains an @-catcode character.
% -----------------------------------------------------------------------------
\makeatletter
\@ifclassloaded{beamer}{%
  \usefonttheme{professionalfonts}
  \setbeamercolor{structure}{fg=primary-blue}
  \setbeamercolor{frametitle}{fg=primary-blue}
  \setbeamercolor{title}{fg=primary-blue}
  \setbeamercolor{subtitle}{fg=primary-gold}
  \setbeamercolor{author}{fg=primary-blue}
  \setbeamercolor{institute}{fg=neutral}
  \setbeamercolor{date}{fg=primary-gold}
  \setbeamercolor{itemize item}{fg=primary-blue}
  \setbeamercolor{itemize subitem}{fg=primary-gold}
  \setbeamercolor{alerted text}{fg=primary-gold}
  \setbeamercolor{block title}{fg=white, bg=primary-blue}
  \setbeamercolor{block body}{bg=light-bg}
  % Semantic block variants: block=definition/concept (blue), exampleblock=
  % worked example (green/positive), alertblock=key takeaway or caution (gold).
  % Keep to <=2 colored boxes per slide across ALL three variants (INV-7).
  \setbeamercolor{block title example}{fg=white, bg=positive}
  \setbeamercolor{block body example}{bg=light-bg}
  \setbeamercolor{block title alerted}{fg=white, bg=primary-gold}
  \setbeamercolor{block body alerted}{bg=light-bg}

  % Minimal footer: course code (left) + page X / N (right), no navigation clutter
  \setbeamertemplate{navigation symbols}{}
  \setbeamertemplate{footline}{%
    \color{neutral}\footnotesize\hspace{0.7em}\@coursecodetext\hfill
    \insertframenumber/\inserttotalframenumber\hspace{0.7em}\vspace{0.3em}}
}{}
\makeatother

% -----------------------------------------------------------------------------
% TikZ setup — libraries the snippet gallery uses
% -----------------------------------------------------------------------------
\usepackage{tikz}
\usetikzlibrary{arrows.meta, positioning, calc, decorations.pathreplacing,
                fit, shapes.geometric, backgrounds}

% Shared TikZ styles — snippets and hand-written diagrams can pick these up
% via `\begin{tikzpicture}[every node/.style={...}]` or by naming specific
% styles (e.g., `\node[dag-node] (X) at (0,0) {X};`).
\tikzset{
  % Node styles for causal DAGs and similar
  dag-node/.style = {
    draw=primary-blue, fill=light-bg, rounded corners=2pt,
    minimum width=1.2cm, minimum height=0.9cm, align=center, font=\small
  },
  decision-node/.style = {
    draw=primary-gold, fill=white, diamond, aspect=1.8,
    minimum width=1.8cm, minimum height=1.2cm, align=center, font=\small,
    inner sep=1pt
  },
  % Edge styles — observed vs counterfactual
  observed-edge/.style = {-{Latex[length=2.5mm]}, thick, draw=jet},
  counterfactual-edge/.style = {-{Latex[length=2.5mm]}, thick, draw=neutral, dashed},
  confound-edge/.style = {-{Latex[length=2.5mm]}, thick, draw=negative, dashed},
  % Data-point styles for plots
  observed-dot/.style = {fill=primary-blue, draw=primary-blue, circle, inner sep=1.2pt},
  counterfactual-dot/.style = {fill=white, draw=neutral, circle, inner sep=1.2pt}
}

% -----------------------------------------------------------------------------
% Convenience macros
% -----------------------------------------------------------------------------
\newcommand{\muted}[1]{\textcolor{neutral}{#1}}
\newcommand{\key}[1]{\textcolor{primary-gold}{\textbf{#1}}}
\newcommand{\good}[1]{\textcolor{positive}{#1}}
\newcommand{\bad}[1]{\textcolor{negative}{#1}}

% Standout frame macro: full-bleed dark background for section transitions.
% Beamer-only (uses \begin{frame}/\setbeamercolor/\usebeamerfont) -- do not
% call from an article-class file (e.g. Notes/<CODE>/*-notes.tex). \muted,
% \key, \good, \bad, and \coursecode above are plain xcolor/gdef and are
% safe to use from either class; verified when Notes/article-class support
% was added.
% Expand: \transitionslide{Your transition title}
\newcommand{\transitionslide}[1]{%
  \begingroup
  \setbeamercolor{background canvas}{bg=primary-blue}
  \begin{frame}[plain]
    \centering
    \vfill
    {\usebeamerfont{title}\color{white}\Large #1\par}
    \vfill
  \end{frame}
  \endgroup
}

% Section divider frame (Beamer-only), an alternative to \transitionslide with a
% darker two-line style: near-black (jet) background, a small white label line
% above a larger gold title, and a thin gold rule along the bottom edge.
% Expand: \sectiondivider{Section 2}{Pointers}
\newcommand{\sectiondivider}[2]{%
  \begingroup
  \setbeamercolor{background canvas}{bg=jet}
  \begin{frame}[plain]
    \vspace*{3.0cm}
    {\color{white}\ttfamily\large #1\par}
    \vspace{0.5cm}
    {\color{primary-gold}\LARGE #2\par}
    \begin{tikzpicture}[remember picture, overlay]
      \draw[primary-gold, line width=2.5pt]
        ($(current page.south west)+(0,0.1)$) -- ($(current page.south east)+(0,0.1)$);
    \end{tikzpicture}
  \end{frame}
  \endgroup
}

% -----------------------------------------------------------------------------
% Channel watermark — "Concepts That Click" (YouTube).
%
% Article-class files (CompetitiveExam Q/A sets, Notes, Assignments) get a
% light diagonal draft watermark on every page plus a centered footer naming
% the channel. Beamer decks get a subtle TikZ background watermark instead
% (the footline already carries the course code). Centralized here so every
% MCQ PDF is branded without per-file edits.
% -----------------------------------------------------------------------------
\makeatletter
\@ifclassloaded{article}{%
  \usepackage{draftwatermark}
  \SetWatermarkText{Concepts That Click}
  \SetWatermarkScale{1}
  \SetWatermarkFontSize{2cm}
  \SetWatermarkLightness{0.86}
  \SetWatermarkAngle{45}
  \usepackage{fancyhdr}
  \pagestyle{fancy}
  \fancyhf{}
  \fancyfoot[C]{\footnotesize\textcolor{neutral}{Concepts That Click~|~YouTube: \href{https://www.youtube.com/@concepts.thatclick}{@concepts.thatclick}}}
  \renewcommand{\headrulewidth}{0pt}
  \renewcommand{\footrulewidth}{0pt}
  \fancypagestyle{plain}{%
    \fancyhf{}%
    \fancyfoot[C]{\footnotesize\textcolor{neutral}{Concepts That Click~|~YouTube: \href{https://www.youtube.com/@concepts.thatclick}{@concepts.thatclick}}}%
    \renewcommand{\headrulewidth}{0pt}%
    \renewcommand{\footrulewidth}{0pt}%
  }
}{}
\@ifclassloaded{beamer}{%
  \setbeamertemplate{background}{%
    \begin{tikzpicture}[remember picture, overlay]
      \node[rotate=45, opacity=0.055, align=center]
        at (current page.center)
        {\Huge\textcolor{neutral}{Concepts That Click}};
    \end{tikzpicture}%
  }
}{}
\makeatother 
\coursecode{JKSSB}
\renewcommand{\thesection}{50.\arabic{section}}
\newtheorem{example}{Example}[section]
\newenvironment{definitionbox}[1]{%
  \par\noindent\textbf{Definition (#1).}\ \itshape
}{\par\vspace{0.3em}}
\title{Data Representation, Number Systems \& Logic --- Detailed Lecture Notes}
\author{JKSSB Computer Awareness}
\date{\today}

\begin{document}
\maketitle

\section{Bases and positional notation}
Digital circuits commonly use two reliably distinguishable states, written
as 0 and 1. A number system specifies a base (also called a radix) and the
digits allowed in each position. The base determines the place values.

\begin{definitionbox}{Number-system base}
The number of distinct digit symbols in a positional number system, with
place values given by powers of that base.
\end{definitionbox}

\begin{center}
\begin{tabular}{llll}
\textbf{Name} & \textbf{Base} & \textbf{Allowed digits} &
\textbf{Example}\\
\hline
Binary & 2 & 0--1 & $1011_2$\\
Octal & 8 & 0--7 & $27_8$\\
Decimal & 10 & 0--9 & $45_{10}$\\
Hexadecimal & 16 & 0--9, A--F & $2F_{16}$\\
\end{tabular}
\end{center}

The subscript identifies the base. In hexadecimal, A represents 10, B 11,
C 12, D 13, E 14, and F 15. A digit that is not allowed in the stated base
makes the numeral invalid: for example, digit 8 cannot appear in an octal
number.

To test whether a numeral is valid, compare every digit with the allowed
range for its base. Binary permits only 0 and 1; octal permits 0 through 7;
decimal permits 0 through 9; hexadecimal permits 0 through 9 and A through
F. A numeral can be converted only after it passes this check. For example,
$128_8$ is not a valid octal numeral because it contains the digit 8, even
though the same characters could be interpreted under a different base.

\begin{definitionbox}{Binary}
A base-2 positional number system using digits 0 and 1.
\end{definitionbox}

\begin{definitionbox}{Octal}
A base-8 positional number system using digits 0 through 7.
\end{definitionbox}

\begin{definitionbox}{Decimal}
A base-10 positional number system using digits 0 through 9.
\end{definitionbox}

\begin{definitionbox}{Hexadecimal}
A base-16 positional number system using digits 0 through 9 and letters A
through F, where A represents 10 and F represents 15.
\end{definitionbox}

Each position contributes its digit multiplied by a power of the base. For
$1011_2$, the rightmost position is $2^0$, the next is $2^1$, then $2^2$,
and then $2^3$:
\[
1011_2=1(2^3)+0(2^2)+1(2^1)+1(2^0)=8+0+2+1=11_{10}.
\]
The rightmost position is the least-significant position because it has the
smallest power.

\begin{center}
\begin{tabular}{ccccc}
\toprule
\textbf{Digit} & 1 & 0 & 1 & 1\\
\textbf{Position from right} & 3 & 2 & 1 & 0\\
\textbf{Place value} & $2^3=8$ & $2^2=4$ & $2^1=2$ & $2^0=1$\\
\textbf{Contribution} & 8 & 0 & 2 & 1\\
\bottomrule
\end{tabular}
\end{center}

The place-value method works in every base. Multiply each digit by the
power assigned to its position, then add the contributions. The rightmost
digit always has power zero; moving one place left increases the exponent
by one. This is why leading zeros do not change a value: each added zero
contributes zero at a new, higher place.

\section{Converting decimal and binary}
To convert a positive decimal integer to binary, divide repeatedly by two.
At each step, keep the remainder, which must be either zero or one. Stop
when the quotient becomes zero, then read the remainders from the last
division upward. Reading them in the order produced gives the wrong
direction.

Use this sequence every time:
\begin{enumerate}
  \item Divide the current integer by 2.
  \item Record the remainder, either 0 or 1.
  \item Replace the current integer with the quotient and repeat until the
    quotient is 0.
  \item Read the recorded remainders from last to first.
  \item Verify the result by adding its binary place-value contributions.
\end{enumerate}

\begin{example}
Convert $45_{10}$ to binary:
\[
\begin{array}{rcl}
45&=&2(22)+1\\
22&=&2(11)+0\\
11&=&2(5)+1\\
5&=&2(2)+1\\
2&=&2(1)+0\\
1&=&2(0)+1
\end{array}
\]
The remainders from top to bottom are $1,0,1,1,0,1$. Reading bottom to top
gives $101101_2$. Check by place value:
$1(32)+0(16)+1(8)+1(4)+0(2)+1(1)=45$, so
$45_{10}=101101_2$.
\end{example}

\begin{example}
Convert $13_{10}$ to binary. Repeated division gives
$13=2(6)+1$, $6=2(3)+0$, $3=2(1)+1$, and $1=2(0)+1$. Reading the
remainders upward gives $1101_2$. Check:
$1(8)+1(4)+0(2)+1(1)=13$.
\end{example}

\section{Binary grouping for octal and hexadecimal}
Three binary bits encode one octal digit because $2^3=8$. Four binary bits
encode one hexadecimal digit because $2^4=16$. To convert a binary numeral,
start at the right and group bits into threes for octal or fours for
hexadecimal. If the leftmost group is short, add leading zeros; leading
zeros do not change the value.

\begin{example}
Convert $1111111_2$ to octal. Group from the right:
$001\ 111\ 111$. The groups are $001_2=1_8$, $111_2=7_8$, and
$111_2=7_8$, so the answer is $177_8$. Check:
$1(8^2)+7(8)+7=64+56+7=127$, the same value as $1111111_2$.
\end{example}

\begin{example}
Convert $1111111_2$ to hexadecimal. Group from the right:
$0111\ 1111$. These groups are $7$ and $15$, and hexadecimal digit 15 is
F. Therefore the answer is $7F_{16}$. Check:
$7(16)+15=112+15=127$.
\end{example}

\begin{example}
Convert $2F_{16}$ to decimal. The hexadecimal digit F is 15:
\[
2F_{16}=2(16^1)+15(16^0)=32+15=47_{10}.
\]
\end{example}

\section{Character codes, Unicode, and decimal digits}
\begin{definitionbox}{ASCII}
A character-encoding standard whose original character set uses 7-bit
codes.
\end{definitionbox}

\begin{definitionbox}{Unicode}
A standard that assigns code points to characters across many writing
systems.
\end{definitionbox}

\begin{definitionbox}{UTF-8}
A variable-length encoding form for Unicode code points.
\end{definitionbox}

\begin{definitionbox}{Binary-Coded Decimal (BCD)}
A representation in which each decimal digit is encoded separately using a
4-bit code.
\end{definitionbox}

\begin{center}
\begin{tabular}{p{0.2\textwidth}p{0.31\textwidth}p{0.37\textwidth}}
\toprule
\textbf{Representation} & \textbf{What it identifies or encodes} & \textbf{Key property}\\
\midrule
ASCII & Characters & Original ASCII uses 7-bit codes\\
Unicode & Characters from many writing systems & Assigns a code point to a character\\
UTF-8 & Unicode code points as encoded bytes & Variable-length encoding\\
BCD & Decimal digits & Uses a separate 4-bit code for each digit\\
\bottomrule
\end{tabular}
\end{center}

ASCII and Unicode concern character representation. Unicode code points
identify characters; UTF-8 is one way to encode those code points as bytes.
The code point and its byte encoding are different things: a character's
Unicode identity does not by itself state how many bytes a particular
encoding uses. UTF-8 is variable-length, so a character is not always stored
using one byte. BCD has a different purpose: it represents each decimal
digit separately using four bits. It is not the same as converting the
entire multi-digit decimal value to binary.

\begin{example}
Represent decimal $59$ in BCD. Encode the digit 5 as $0101$ and the digit 9
as $1001$, giving $0101\ 1001$. Ordinary binary for the numeric value 59 is
$111011_2$. These bit strings differ because BCD preserves the two decimal
digits separately, while ordinary binary represents the value 59 as one
base-2 numeral.
\end{example}

The grouping explains why BCD and ordinary binary use different numbers of
bits. BCD uses one 4-bit group for each decimal digit, so a two-digit
decimal value uses two groups. Ordinary binary instead places the entire
quantity in powers of 2. When an item asks for BCD, split the decimal
number into its digits before encoding; do not convert the whole number
directly to binary.

\section{RGB colour representation}
\begin{definitionbox}{RGB}
A colour model representing colour through red, green, and blue channel
values.
\end{definitionbox}

In a common 24-bit RGB representation, each of the three channels uses 8
bits. Therefore the three channels require $8+8+8=24$ bits per pixel.
Eight bits can encode $2^8=256$ possible values per channel, from 0 through
255. The total number of possible channel combinations in that representation
is $256^3=16{,}777{,}216$. This does not mean every image format or display
uses exactly 24-bit RGB; other bit depths, alpha channels, palettes, or
colour spaces are possible.

The channel count and the bit count are separate quantities. RGB has three
channels: red, green, and blue. In the common representation described
above, each channel has 8 bits, so one pixel uses 24 bits total. A value
within one channel ranges from 0 to 255 because an 8-bit field has
$2^8=256$ possible patterns. To count all channel combinations, multiply
the number of choices for each independent channel:
\[
  256\times256\times256=256^3=16{,}777{,}216.
\]
This count follows only under the stated 8-bits-per-channel assumption.

\begin{example}
In common 24-bit RGB, the red, green, and blue channels each use 8 bits.
Adding the channel widths gives $8+8+8=24$ bits per pixel.
\end{example}

\section{Logic gates and truth tables}
\begin{definitionbox}{AND gate}
A logic operation that outputs 1 only when every input is 1.
\end{definitionbox}

\begin{definitionbox}{OR gate}
A logic operation that outputs 1 when at least one input is 1.
\end{definitionbox}

\begin{definitionbox}{XOR gate}
A logic operation that outputs 1 when its two inputs differ.
\end{definitionbox}

\begin{definitionbox}{NAND gate}
The negation of AND: it outputs 0 only when every input is 1.
\end{definitionbox}

\begin{definitionbox}{NOT gate}
A one-input logic operation that reverses its input: 0 becomes 1 and 1
becomes 0.
\end{definitionbox}

\begin{center}
\begin{tabular}{cc|cccc|c}
\toprule
$A$ & $B$ & AND & OR & XOR & NAND & NOT $A$\\
\midrule
0 & 0 & 0 & 0 & 0 & 1 & 1\\
0 & 1 & 0 & 1 & 1 & 1 & 1\\
1 & 0 & 0 & 1 & 1 & 1 & 0\\
1 & 1 & 1 & 1 & 0 & 0 & 0\\
\bottomrule
\end{tabular}
\end{center}

Read the truth table from the rule, not by memorizing an unlabelled row.
AND requires both inputs to be 1. OR needs at least one 1. XOR checks
whether the inputs differ. NAND first evaluates AND and then reverses that
output. NOT has only one input and reverses it.

For two-input gates there are exactly four possible input pairs:
$00$, $01$, $10$, and $11$. Evaluate every pair using the gate's rule:
\begin{itemize}
  \item AND is 1 only on the $11$ row.
  \item OR is 0 only on the $00$ row.
  \item XOR is 1 on $01$ and $10$, where the inputs differ.
  \item NAND is the opposite of AND on every row.
  \item NOT does not use $B$; it reverses $A$ in each row.
\end{itemize}
OR and XOR are easy to confuse. For input $A=1$, $B=1$, OR is 1 because
at least one input is 1, but XOR is 0 because the two inputs are equal.
NAND is checked only after evaluating AND and then taking the opposite
result.

\begin{example}
For $A=1$ and $B=0$, AND is 0 because both inputs are not 1; OR is 1
because one input is 1; XOR is 1 because the values differ; and NAND is 1
because it reverses the AND result. NOT($A$) is 0.
\end{example}

\begin{example}
For $A=1$ and $B=1$, AND is 1 and OR is 1. XOR is 0 because the inputs do
not differ. NAND is the opposite of AND, so it is 0. NOT($A$) is 0.
\end{example}

\subsection*{Exercises}
\begin{enumerate}[label=\arabic*.]
  \item Convert $29_{10}$ to binary and verify the result by place value.
  \item Convert $101101_2$ to octal and hexadecimal.
  \item Explain the difference between BCD for decimal 25 and ordinary
    binary for the value 25.
  \item In common 24-bit RGB, how many bits are assigned to each channel?
    Is that rule universal for all images?
  \item For $A=0$, $B=1$ and for $A=1$, $B=1$, find AND, OR, XOR, NAND,
    and NOT($A$).
\end{enumerate}

\subsection*{Solutions}
\begingroup\small
\begin{enumerate}[label=\arabic*.]
  \item Repeated division gives $29=2(14)+1$, $14=2(7)+0$,
    $7=2(3)+1$, $3=2(1)+1$, and $1=2(0)+1$. Reading upward gives
    $11101_2$. Check: $16+8+4+0+1=29$.
  \item Group from the right: $101\ 101$, so the octal result is $55_8$.
    Pad and group by four: $0010\ 1101$, so the hexadecimal result is
    $2D_{16}$.
  \item BCD for 25 is $0010\ 0101$, encoding 2 and 5 separately.
    Ordinary binary for the value 25 is $11001_2$.
  \item Eight bits per channel, since $8+8+8=24$. No; this describes a
    common 24-bit representation, not every image or display.
  \item For $A=0$, $B=1$: AND=0, OR=1, XOR=1, NAND=1, and NOT($A$)=1.
    For $A=1$, $B=1$: AND=1, OR=1, XOR=0, NAND=0, and NOT($A$)=0.
\end{enumerate}
\endgroup
\end{document}
